\subsection{Primitive relative spectral dynamics}

The carrier theorem is not intended to float above an unspecified microscopic dynamics. The
starting point is therefore the one-sided boundary datum from which the primitive layer, the
derived closed datum, and the canonical admissibility projector are reconstructed before the later
closure theorems are invoked.

\begin{definition}[One-sided boundary admissibility datum]
\label{def:reduced-theorem-datum}
The one-sided boundary admissibility datum is
\[
\Tbound:=(\mathcal A_+,\mathcal H_+,D_+,J,\Gamma,B_\Sigma)
\]
with the following built-in admissibility conditions:
\begin{enumerate}[label=(\roman*), leftmargin=1.5em, topsep=3pt, itemsep=2pt]
    \item $Z(\mathcal A_+)$ is spectrally regular and its Gelfand spectrum is a smooth compact oriented manifold with boundary $M_+$;
    \item $B_\Sigma$ is a self-adjoint tangential boundary operator on $\Sigma:=\partial M_+$ compatible with the boundary restrictions of $J$ and $\Gamma$;
    \item in a collar neighborhood of $\Sigma$, the Dirac operator has canonical product form
    \[
    D_+=\gamma_n(\partial_n+B_\Sigma);
    \]
    \item the boundary collar is reflection-regular so that canonical doubling, the Calder\'on projector, and the admissibility complex are well posed on the same datum.
\end{enumerate}
\end{definition}

\begin{remark}[Almost-commutative realizations]
\label{ass:almost-commutative-factorization}
When a factorized realization is useful, one may choose
\[
\mathcal H\cong L^2(\widetilde M,S)\otimes \Hint,
\]
with $S\to\widetilde M$ the spinor bundle on the oriented double cover and $\Hint$ a finite
internal Hilbert space. This is an analytic realization choice rather than part of the boundary
datum.
\end{remark}

\begin{theorem}[Doubling, deck involution, and boundary polarization]
\label{thm:primitive-datum-from-reduced-datum}
\label{thm:primitive-scaffold-from-reduced-datum}
Let $\Tbound$ be as above. Let
\[
\widetilde M:=M_+\cup_{\Sigma} M_+^{\mathrm{op}}
\]
be the doubled manifold and let
\[
\taudbl:\widetilde M\to\widetilde M
\]
be its deck involution. Let $P_C$ be the Calder\'on projector of the APS-realized boundary problem
for $D_+$ on $\Sigma$, and define the boundary polarization operator
\[
\iC:=2P_C-1
\]
on the Cauchy-data Hilbert space $\mathcal H_{\partial}$.

Then:
\begin{enumerate}[label=(\roman*), leftmargin=1.8em, topsep=3pt, itemsep=2pt]
    \item the geometric objects are
    \[
    M=\widetilde M/\taudbl,
    \qquad
    \Seam=\mathrm{Fix}(\taudbl),
    \qquad
    \Xbulk=\widetilde M\setminus\Seam;
    \]
    \item the doubled operator and its even and odd parts are
    \[
    D=D_+\cup_{\Sigma}D_+^{\mathrm{op}},
    \qquad
    D_{\mathrm{ref}}:=\frac12\bigl(D+\taudbl^{\ast}D\taudbl^{\ast}\bigr),
    \qquad
    D_{\mathrm{rel}}:=\frac12\bigl(D-\taudbl^{\ast}D\taudbl^{\ast}\bigr);
    \]
    \item on boundary Cauchy data the induced reflection $R_{\taudbl}$ satisfies
    \[
    \operatorname{Ran}(P_C)=\ker(R_{\taudbl}-1),
    \qquad
    \ker(P_C)=\ker(R_{\taudbl}+1),
    \qquad
    \iC=R_{\taudbl}.
    \]
\end{enumerate}
Hence the primitive scaffold
\[
\Tprim^{\mathrm{kin}}
:=
(\Xbulk,\widetilde M,M,\Seam,\taudbl,\iC,D_{\mathrm{ref}},D_{\mathrm{rel}},\chiseed)
\]
is derived, not assumed. The derived doubled datum may be recorded as
\[
\Tmin^{\mathrm{cl}}:=(\mathcal A,\mathcal H,D,J,\Gamma,\taudbl,\iC),
\]
If one further chooses an almost-commutative factorization
\[
\mathcal H_{\mathrm{prim}}\cong L^2(\Xbulk,S)\otimes \Hint,
\]
then $S$ and $\Hint$ refer only to that chosen realization and are not additional
reconstruction outputs.
\end{theorem}

\begin{proof}
Reflection regularity on an exact collar produces the doubled manifold together with the deck
involution $\taudbl$. The APS-realized boundary problem has a Calder\'on projector $P_C$ whose
range is the Cauchy-data space of interior solutions. On the doubled collar the deck reflection
acts on Cauchy data by $R_{\taudbl}$, and the solution space splits into the $\pm1$ eigenspaces of
this reflection. Therefore
\[
P_C=\frac{1+R_{\taudbl}}{2},
\qquad
\iC=2P_C-1=R_{\taudbl}.
\]
The formulas for $D_{\mathrm{ref}}$ and $D_{\mathrm{rel}}$ are the even and odd parts of $D$ under
the induced involution.
\end{proof}

\begin{corollary}[No confusion among the three involutions]
Geometric reflection is carried by $\taudbl$, boundary polarization by $\iC$, and the finite
carrier split by the restricted involution
\[
\epscar:=\iC|_{\mathcal H_{\mathrm{car}}}.
\]
No later theorem may identify these three operators literally; only the displayed restriction and
transport maps are allowed.
\end{corollary}

The canonical admissibility projector is built in two transparent stages. The primitive stage
uses only data already available from $\Tbound$ through the derived boundary polarization $\iC$,
namely the seam-even selector on Cauchy data, an APS realization that makes the geometric anchor
self-adjoint Fredholm, and the relative operator $D_{\mathrm{rel}}$. The full stage upgrades this
primitive projector by the color and determinant projectors that become available only after the
carrier theorem of \Cref{sec:minimal-carrier-proof} identifies the color factor and after the
compact bosonic index of \Cref{sec:closure-theorems} fixes the determinant involution. The full
complex and the upgrade theorem appear once these later inputs are in place; here we close the
primitive stage.

\begin{definition}[Primitive admissibility complex and primitive projector]
\label{def:admissibility-complex}
\label{def:primitive-admissibility-complex}
Let the primitive Hilbert space carry only the seam factor and the gap sector accessible from
$D_{\mathrm{rel}}$ alone:
\[
\mathcal H_{\mathrm{prim}}=\mathcal H_{\Seam}\otimes\mathcal H_t,
\qquad
\mathcal H_{\mathrm{prim}}^{\mathrm{ext}}
:=
\mathcal H_{\mathrm{prim}}\otimes\Lambda(\eta_\Sigma,\eta_g),
\]
with primitive seam selector
\[
P_{\Seam,+}:=\frac{\mathbf 1+\iC}{2}.
\]
Fix $\epsilon_{\mathrm{adm}}>0$ and let $D_{\mathrm{prim}}$ be the APS-realized restriction of
the primitive sector operator to $\mathrm{Ran}(P_{\Seam,+})$, controlled by $D_{\mathrm{rel}}$
relative to the geometric anchor. Define the primitive admissibility penalties
\[
K_\Sigma:=\sqrt{\lambda_\Sigma}\,\frac{\mathbf 1-\iC}{2},
\qquad
K_g:=\sqrt{\lambda_{\mathrm{gap}}}\,
\mathbf 1_{[0,\epsilon_{\mathrm{adm}}^2/4)}\!
\bigl(P_{\Seam,+}D_{\mathrm{prim}}^\dagger D_{\mathrm{prim}}P_{\Seam,+}\bigr),
\]
the primitive admissibility differential
\[
Q_{\mathrm{prim}}
:=
\eta_\Sigma^\dagger K_\Sigma
+\eta_g^\dagger K_g,
\]
and the primitive admissibility Laplacian
\[
\Delta_{\mathrm{prim}}:=\{Q_{\mathrm{prim}},Q_{\mathrm{prim}}^\dagger\}.
\]
The primitive admissibility projector is the Hodge projector
\[
\Pprim
:=
\Pi_{\ker\Delta_{\mathrm{prim}}}.
\]
\end{definition}

\begin{lemma}[Primitive nilpotency and primitive Hodge projector]
\label{lem:admissibility-nilpotency}
\label{lem:hodge-admissibility-projector}
\label{lem:primitive-admissibility-nilpotency}
On the primitive scaffold reconstructed from $\Tbound$,
\[
Q_{\mathrm{prim}}^2=0,
\qquad
\Delta_{\mathrm{prim}}
=
K_\Sigma^2+K_g^2,
\qquad
\Pprim
=
\operatorname*{s-lim}_{\beta\to\infty}e^{-\beta(K_\Sigma^2+K_g^2)}.
\]
In particular $\Pprim$ depends only on $\iC$, the APS realization, and $D_{\mathrm{rel}}$, and
is therefore reconstructed from the closed datum at the primitive level. Equivalently,
\[
\Pprim=P_{\Seam,+}P_{\mathrm{prim,gap}},
\qquad
P_{\mathrm{prim,gap}}
:=
\mathbf 1_{[\epsilon_{\mathrm{adm}}^2/4,\infty)}
\!\bigl(P_{\Seam,+}D_{\mathrm{prim}}^\dagger D_{\mathrm{prim}}P_{\Seam,+}\bigr).
\]
\end{lemma}

\begin{proof}[Proof sketch]
The auxiliary generators $\eta_\Sigma^\dagger$ and $\eta_g^\dagger$ anticommute, so the square
of $Q_{\mathrm{prim}}$ reduces to commutators of $K_\Sigma$ and $K_g$. Both penalties act on
commuting sectors of $\mathcal H_{\mathrm{prim}}$, hence the mixed terms vanish and
$\Delta_{\mathrm{prim}}=K_\Sigma^2+K_g^2$. The zero-temperature limit projects onto the
intersection of seam-even and gapped admissible vectors, which is the factorized form recorded
above.
\end{proof}

\begin{definition}[Carrier polarization on the finite block]
\label{def:carrier-polarization-involution}
Let $\mathcal H_{\mathrm{car}}$ denote the finite-dimensional carrier block selected after
primitive admissibility compression. Define
\[
\epscar:=\iC|_{\mathcal H_{\mathrm{car}}},
\qquad
E_-:=\ker(\epscar+1),
\qquad
E_+:=\ker(\epscar-1).
\]
Then the finite carrier decomposition is
\[
E=E_-\oplus E_+.
\]
\end{definition}

\begin{remark}[Forward reference to the full admissibility complex]
\label{rem:forward-Padm}
The full admissibility projector is obtained by upgrading $\Pprim$ with the color center
projector $P_{\mathrm{sing}}$ and the determinant involution projector $P_\Theta$, which are
available only after the carrier theorem of \Cref{sec:minimal-carrier-proof} identifies the
color factor $\mathcal H_c$ and after the compact bosonic index of \Cref{sec:closure-theorems}
fixes the determinant involution $Q_\Theta$. The full complex and the explicit upgrade are
stated in \Cref{def:full-admissibility-complex,thm:padm-upgrade}; nothing in the present
primitive subsection depends on those inputs.
\end{remark}

\begin{definition}[Primitive Euclidean action]
For any declared spectral test profile $f$, the primitive Euclidean action is the relative
spectral seed
\[
S_{\mathrm{prim}}[\chiseed]
:=
\tr f\!\left(\frac{D_{\mathrm{rel}}}{\chiseed}\right)
-
\tr f\!\left(\frac{D_{\mathrm{ref}}}{\chiseed}\right)
+\frac{i\pi}{2}\Delta\eta_\Sigma.
\]
Gauge fixing, ghosts, and carrier-compatible bosonic fluctuations are introduced only after the
hard carrier theorem and the canonical admissibility projector are in place. Whenever convenient,
the compressed expression $\Padm A\Padm$ may be used as shorthand for matrix elements on the
physical admissible sector, but it is not the definition of a local operator net.
\end{definition}

\begin{definition}[Physical admissible local net]
\label{def:physical-admissible-local-net}
Let $\mathfrak A_E(\mathcal O)$ be the Euclidean source algebra generated by local field insertions
with test-function support in $\mathcal O$. After the later Osterwalder--Schrader reconstruction,
let $\mathfrak A_M(\mathcal O)$ denote the corresponding Minkowski local algebra and let
$\mathcal H_{\mathrm{OS}}$ be the reconstructed Hilbert space. Define the physical admissible
Hilbert space by
\[
\mathcal H_{\mathrm{adm}}:=\operatorname{Ran}(\Padm)\subset \mathcal H_{\mathrm{OS}}.
\]
Define the physical admissible local net by
\[
\mathfrak A_{\mathrm{adm}}(\mathcal O)
:=
\left\{
A|_{\mathcal H_{\mathrm{adm}}}\;:\;
A\in\mathfrak A_M(\mathcal O),\;
A\mathcal H_{\mathrm{adm}}\subset \mathcal H_{\mathrm{adm}},\;
A^\ast\mathcal H_{\mathrm{adm}}\subset \mathcal H_{\mathrm{adm}}
\right\}.
\]
Thus $\Padm$ acts as the selector of the physical state space, while the local net is defined by
restriction to the invariant admissible subspace rather than by compression with a global
projector.
\end{definition}

\begin{remark}[Bernstein realizations as analytic choices]
\label{ass:bernstein-positivity}
Whenever a Bernstein representation of a spectral test profile is used, it is treated only as an
analytic realization of a transfer channel already fixed by the closed datum and the canonical
projector. Positivity or reflection stability therefore does not enter the main text as a
primitive source of admissibility.
\end{remark}

\begin{theorem}[Primitive self-adjoint realization without extra analytic assumption]
\label{thm:well-posed-primitive-dynamics}
\label{thm:primitive-selfadjoint-realization}
Let $M_\Lambda$ be a finite cutoff of the primitive geometric anchor, so that $M_\Lambda$ is
compact with boundary. Let $D_{\mathrm{geo}}$ be the formally self-adjoint Dirac-type anchor and
let $A_\partial$ be an adapted boundary operator. Choose a self-adjoint $D_{\mathrm{geo}}$-elliptic
boundary condition $B_{\mathrm{sa}}$ of modified APS type:
\[
B_{\mathrm{sa}}=
\begin{cases}
H^{1/2}_{(-\infty,0)}(A_\partial), & \ker A_\partial=0,\\[2mm]
H^{1/2}_{(-\infty,0)}(A_\partial)\oplus L, & \ker A_\partial\neq0,
\end{cases}
\]
where $L\subset\ker A_\partial$ is chosen as in the normal form for self-adjoint $D$-elliptic
boundary conditions. Then the closed realization $D_{\mathrm{geo},B_{\mathrm{sa}}}$ is self-adjoint
Fredholm with discrete spectrum. If $D_{\mathrm{rel}}$ is a symmetric differential perturbation of
order at most one with bounded coefficients on $M_\Lambda$, then $D_{\mathrm{rel}}$ is
infinitesimally $D_{\mathrm{geo},B_{\mathrm{sa}}}$-bounded and therefore
\[
D_{\mathrm{geo},B_{\mathrm{sa}}}+D_{\mathrm{rel}}
\]
is self-adjoint on the same domain. In particular the finite-cutoff relative spectral seed is well
defined without an additional analytic existence hypothesis.
\end{theorem}

\begin{proof}
The existence of self-adjoint $D$-elliptic boundary conditions for formally self-adjoint Dirac-type
operators is standard. In the zero-kernel case APS is already self-adjoint. If the adapted
boundary operator has nontrivial kernel, one passes to a self-adjoint modified APS realization or,
equivalently, to
\[
H^{1/2}_{(-\infty,0)}(A_\partial)\oplus L,
\]
with $L\subset\ker A_\partial$ chosen as in the normal form theorem for self-adjoint $D$-elliptic
boundary conditions. The resulting realization is essentially self-adjoint on smooth sections
satisfying the boundary condition. Since $M_\Lambda$ is compact, coercivity at infinity is
automatic and the realization is Fredholm with discrete spectrum.

For the perturbation, standard first-order elliptic regularity on a compact manifold with
self-adjoint $D$-elliptic boundary condition gives a graph norm estimate
\[
\|\psi\|_{H^1}\le C\bigl(\|D_{\mathrm{geo},B_{\mathrm{sa}}}\psi\|+\|\psi\|\bigr).
\]
Because $D_{\mathrm{rel}}$ has order at most one and bounded coefficients,
\[
\|D_{\mathrm{rel}}\psi\|
\le
C_1\|\psi\|_{H^1}+C_0\|\psi\|
\le
\varepsilon\|D_{\mathrm{geo},B_{\mathrm{sa}}}\psi\|+C_\varepsilon\|\psi\|
\]
for every $\varepsilon>0$. Hence $D_{\mathrm{rel}}$ is $D_{\mathrm{geo},B_{\mathrm{sa}}}$-bounded
with relative bound zero. The Kato--Rellich theorem therefore yields self-adjointness of
$D_{\mathrm{geo},B_{\mathrm{sa}}}+D_{\mathrm{rel}}$ on the same domain.
\end{proof}

\begin{remark}[APS roles on the admissible branch]
\label{rem:aps-roles}
At the primitive level $\APSdomprov$ is the provisional elliptic boundary choice for the
geometric anchor: an analytic starting representative chosen so that
\Cref{thm:well-posed-primitive-dynamics} applies. After the carrier theorem and the master
closure roadmap of \Cref{sec:closure-theorems} the Calder\'on projector of the doubled operator
$D_b$ canonically promotes $\APSdomprov$ to the closed-branch boundary domain $\APSdom$ in the
same admissible class. The family-pairing data $\APSfam$ are then induced canonically and remain
a derived consequence rather than a second primitive input package. The explicit upgrade
statement is recorded as \Cref{thm:apsdom-upgrade} in \Cref{sec:closure-theorems}.
\end{remark}

\begin{figure}[t]
\centering
\scriptsize
\begin{tikzpicture}[
    root/.style={draw=black!65, fill=gray!10, rounded corners, text width=3.3cm, minimum height=0.9cm, align=center},
    prov/.style={draw=blue!60!black, fill=blue!6, rounded corners, text width=3.15cm, minimum height=1.45cm, align=center},
    dom/.style={draw=blue!70!black, fill=blue!10, rounded corners, very thick, text width=3.15cm, minimum height=1.45cm, align=center},
    fam/.style={draw=orange!75!black, fill=orange!8, rounded corners, text width=3.15cm, minimum height=1.45cm, align=center},
    note/.style={font=\footnotesize\itshape, text=black!70, align=center}
]
\node[root] (aps) at (0,0.45) {APS roles in the present manuscript};
\node[prov] (prov) at (-4.05,-1.95) {\textbf{$\APSdomprov$}\\provisional APS realization\\(analytic starting choice)};
\node[dom]  (dom)  at (0,-1.95)    {\textbf{$\APSdom$}\\closed-branch APS domain\\via Calder\'on upgrade};
\node[fam]  (fam)  at (4.05,-1.95)  {\textbf{$\APSfam$}\\family pairing induced\\by canonical doubling};
\draw[->, thick] (aps.south west) -- ($(prov.north)+(0,0.05)$);
\draw[->, thick] (aps.south) -- ($(dom.north)+(0,0.05)$);
\draw[->, thick] (aps.south east) -- ($(fam.north)+(0,0.05)$);
\draw[->, thick, blue!70!black] ($(prov.north east)+(0.08,0.16)$)
    .. controls +(0.7,0.75) and +(-0.7,0.75) ..
    node[above, fill=white, inner sep=1.5pt, font=\scriptsize\itshape, text=blue!70!black]
    {Calder\'on upgrade}
    ($(dom.north west)+(-0.08,0.16)$);
\draw[->, dashed, black!45, line width=0.8pt] ($(dom.north east)+(0.08,0.14)$)
    .. controls +(0.7,0.65) and +(-0.7,0.65) ..
    node[above, fill=white, inner sep=1.5pt, font=\scriptsize\itshape, text=black!65]
    {induced}
    ($(fam.north west)+(-0.08,0.14)$);
\node[note] at (0,-3.2) {Provisional choice $\to$ Calder\'on upgrade $\to$ derived family pairing.};
\end{tikzpicture}
\caption{APS organization with Calder\'on upgrade and induced family pairing.}
\label{fig:aps-split}
\end{figure}

\begin{corollary}[Carrier-compatible Lorentzian dynamics after closure]
Once the later carrier-compatible fluctuation package and the admissibility closure hypotheses of \Cref{sec:closure-theorems} are imposed, the resulting admissible Euclidean theory reconstructs a self-adjoint Lorentzian Hamiltonian $H_{\mathrm{adm}}$ and a one-parameter automorphism group
\[
\alpha_t(O)=e^{itH_{\mathrm{adm}}}Oe^{-itH_{\mathrm{adm}}}.
\]
\end{corollary}

